This narrative portrays escalating tension between the Navajo leader Black Horse and the Navajo Agent Dana Shipley. Agent Shipley is sent to Round Rock to execute the law passed in 1887 according to which all children needed to attend school. The fight starts with verbal exchanges and intensifies into physical confrontation, until Black Horse challenges the authority. Amid the chaos, notable figures, including Sucker and Chee Dodge, strive to temper Black Horse's aggression and mediate peace. When soldiers arrive, both sides brace for conflict, yet the appearance of a War Chief (Officer) shifts the course, leading Black Horse to lay down his weapon and extend a handshake to signal the end of hostilities.

This version focuses on the immediacy of the conflict as a fierce struggle for identity and collective resistance and is therefore more politically charged than the other versions of the Trouble at Round Rock. Left-Handed Mexican Clansman, a witness to this historical event, presents Black Horse as a heroic figure who battles against the colonial forces with emotional stakes involved in the struggle over education. Concluding with \textit{Íídą́ą́ t’ah chąąmą’ii nishłį́įgo, éí t’éiyá biniiyé ’áadi ’atah nahashłe’} (`At that time, I was just a young punk -- that’s probably why I was just there with the crowd for fun'), the author suggests that his perspective on the event may have evolved since then.